\documentclass[aspectratio=169,10pt,t]{beamer}
\usepackage{ac1}
\aulainfo{09}{Laboratório de Problemas Reais}{Validação de dados, médias e regras de negócio}

\begin{document}
\capa

\begin{frame}{Plano do laboratório}
  \begin{itemize}
    \item Aula prática: o computador está aberto, o livro está fechado
    \item Problemas com cara de \textbf{engenharia}: validar entradas, calcular médias, aplicar regras
    \item Método: \alert{entender \(\to\) planejar \(\to\) codificar \(\to\) testar}
    \item Em duplas; ao fim de cada problema, compare com o site do livro
  \end{itemize}
  \begin{block}{Revisão rápida: ferramentas já disponíveis}
    \texttt{input}/\texttt{print}, tipos (\texttt{int}, \texttt{float}, \texttt{str}), \texttt{if}/\texttt{elif}/\texttt{else}, \texttt{while}, \texttt{for}/\texttt{range}.
  \end{block}
\end{frame}

\begin{frame}{Método de resolução em 4 passos}
  \begin{enumerate}
    \item \textbf{Entender:} entradas, saídas, restrições. Qual é um exemplo com resposta conhecida?
    \item \textbf{Planejar:} esboço em pseudocódigo ou fluxograma
    \item \textbf{Codificar:} um pedaço por vez; execute a cada passo
    \item \textbf{Testar:} casos típicos, limites e entradas inválidas
  \end{enumerate}
  \vspace{1ex}
  \begin{alertblock}{Erro comum}
    Começar digitando código sem ter claro o que o programa deve fazer. Capítulo 9 do livro: \emph{Pense antes de programar}.
  \end{alertblock}
\end{frame}

\begin{frame}{Problema 1: média com validação}
  \begin{exampleblock}{Enunciado}
    Leia as notas de \(n\) provas de um aluno (\(0 \le \text{nota} \le 10\)). Rejeite notas inválidas pedindo novamente. Ao final, mostre a média e a situação:
    \begin{itemize}\small
      \item média \(\ge 7\): aprovado
      \item \(4 \le\) média \(< 7\): prova final
      \item média \(< 4\): reprovado
    \end{itemize}
  \end{exampleblock}
  \textbf{Perguntas guia}
  \begin{itemize}
    \item Como pedir \emph{novamente} até a nota ser válida? \hfill (\texttt{while})
    \item Como somar \(n\) valores? \hfill (\texttt{for} + acumulador)
    \item Quais casos de teste? \hfill (notas 0, 10, 3.9, 4.0, 7.0)
  \end{itemize}
\end{frame}

\begin{frame}[fragile]{Problema 1: solução comentada}
\begin{lstlisting}
n = int(input("Número de provas: "))
soma = 0
for i in range(1, n + 1):
    nota = float(input(f"Nota da prova {i}: "))
    while nota < 0 or nota > 10:          # validação
        print("Nota inválida! Use valores entre 0 e 10.")
        nota = float(input(f"Nota da prova {i}: "))
    soma += nota
media = soma / n
print(f"Média: {media:.2f}")
if media >= 7:
    print("Aprovado")
elif media >= 4:
    print("Prova final")
else:
    print("Reprovado")
\end{lstlisting}
\end{frame}

\begin{frame}{Problema 2: tabela de conversão}
  \begin{exampleblock}{Enunciado (Exercício 3.14)}
    Escreva um programa que imprima uma tabela de temperaturas em graus \textbf{Celsius} de \(-20\) a \(40\), de 10 em 10, com a temperatura equivalente em \textbf{Fahrenheit}.
  \end{exampleblock}
  \[ F = \frac{9}{5}\,C + 32 \]
  \textbf{Perguntas guia}
  \begin{itemize}
    \item Qual laço? Qual \texttt{range}? \hfill (\texttt{range(-20, 41, 10)})
    \item Como alinhar colunas? \hfill (formatação \texttt{\{:>8.1f\}})
  \end{itemize}
\end{frame}

\begin{frame}[fragile]{Problema 2: solução}
\begin{lstlisting}
print(f"{'Celsius':>8}{'Fahrenheit':>12}")
for c in range(-20, 41, 10):
    f = 9 / 5 * c + 32
    print(f"{c:>8}{f:>12.1f}")
\end{lstlisting}
\begin{lstlisting}[style=saida]
 Celsius  Fahrenheit
     -20        -4.0
     -10        14.0
       0        32.0
      10        50.0
...
\end{lstlisting}
  \small Variante: o usuário informa o início, o fim e o passo (e o programa valida que \(\text{passo}>0\)).
\end{frame}

\begin{frame}{Problema 3: regra de desconto}
  \begin{exampleblock}{Enunciado}
    Uma loja concede desconto sobre o valor da compra:
    \begin{itemize}\small
      \item compras até R\$\,100: sem desconto
      \item de R\$\,100 a R\$\,500 (exclusive): 5\%
      \item R\$\,500 ou mais: 10\%; e \textbf{mais} 2\% se o cliente for ``fidelidade''
    \end{itemize}
    Leia valor e se há cartão fidelidade (\texttt{s}/\texttt{n}); mostre o valor final.
  \end{exampleblock}
  \textbf{Pontos de atenção}
  \begin{itemize}
    \item Aceitar \texttt{S} e \texttt{s}; rejeitar outras respostas (validação)
    \item Descontos \emph{somam} ou \emph{compõem}? Defina a regra antes de programar!
  \end{itemize}
\end{frame}

\begin{frame}[fragile]{Problema 3: solução}
\begin{lstlisting}
valor = float(input("Valor da compra: R$ "))
resp = input("Cliente fidelidade (s/n)? ").lower()
while resp not in ("s", "n"):
    resp = input("Responda s ou n: ").lower()

if valor >= 500:
    desc = 0.10
elif valor >= 100:
    desc = 0.05
else:
    desc = 0.0
if valor >= 500 and resp == "s":
    desc += 0.02
final = valor * (1 - desc)
print(f"Desconto de {desc*100:.0f}% -> R$ {final:.2f}")
\end{lstlisting}
\end{frame}

\begin{frame}{Problema 4: série histórica simples}
  \begin{exampleblock}{Enunciado}
    Leia as temperaturas diárias (em \textdegree C) de uma semana (7 valores). Mostre a \textbf{média}, a \textbf{maior} e a \textbf{menor} temperatura e quantos dias ficaram acima da média.
  \end{exampleblock}
  \begin{itemize}
    \item Dá para resolver \alert{sem listas}? Maior e menor com variáveis auxiliares, mas ``dias acima da média'' exige guardar os valores
    \item Esse é o motivo de existirem \textbf{listas} (Aula 16 em diante)
  \end{itemize}
  \begin{block}{Desafio extra para quem terminar}
    Resolva sem listas as três primeiras partes e anote o que não foi possível. Guarde sua solução: voltaremos a ela na Aula 17.
  \end{block}
\end{frame}

\begin{frame}{Depuração: erros comuns observados}
  \centering\small
  \begin{tabular}{p{4.2cm}p{7.2cm}}
    \toprule
    Sintoma & Causa provável\\
    \midrule
    \texttt{ValueError} no \texttt{int(input())} & usuário digitou texto ou deixou em branco\\
    Média errada & esqueceu de converter com \texttt{float}; dividiu por constante errada\\
    Laço infinito & variável da condição não é atualizada dentro do laço\\
    Última nota some / sobra uma & \texttt{range(n)} versus \texttt{range(1, n+1)}\\
    \texttt{IndentationError} & mistura de tabulação e espaços\\
    Desconto aplicado errado & ordem dos \texttt{if}/\texttt{elif} invertida\\
    \bottomrule
  \end{tabular}
\end{frame}

\exercicios{3}{%
  \begin{itemize}\small
    \item Faixas e regras: \textbf{3.7, 3.8, 3.9, 3.10}
    \item Médias e repetição: \textbf{3.11, 3.12, 3.13, 3.14}
    \item Padrões: \textbf{3.15}
  \end{itemize}
  \vspace{0.5ex}
  \footnotesize Compare suas soluções com as do site depois de tentar sozinho.}

\resumo{
  \item Validar a entrada com \texttt{while} é padrão em programas reais
  \item Primeiro as regras no papel, depois o código
  \item Teste os valores de fronteira (0, 4, 7, 10, \ldots)
  \item Próxima aula: subalgoritmos (funções)
}

\end{document}
